\section{问题简介}

图像分类是计算机视觉中的基础任务，而在真实应用中，研究者经常会遇到一个很实际的问题：目标数据集规模较小，但模型结构已经较成熟，此时究竟应该从头训练，还是直接利用公开预训练权重进行微调。这个问题在课程实验中也很有代表性，因为它不仅关系到最终分类精度，还会影响训练时间、参数收敛速度以及调参难度。

本次实验围绕这一问题展开。本文选择两种典型的卷积神经网络结构，即 ResNeXt50\_32x4d 和 DenseNet121，在 Beans 三分类数据集上分别进行从头训练和微调，共形成四组对照实验。这样的设置有两个考虑：一方面，ResNeXt 与 DenseNet 在特征提取方式上存在明显差异，前者强调 grouped transformation 与 cardinality，后者强调特征复用和密集连接；另一方面，小规模农业病害数据集也比较适合观察迁移学习在有限样本条件下的效果。

本实验希望回答三个问题。第一，在小样本图像分类任务中，微调是否一定优于从头训练。第二，不同 CNN 结构在同一数据集上的表现是否一致。第三，在 CPU 条件下完成训练时，哪一种训练策略更稳定、更容易复现。从课程作业的角度看，这三个问题既有理论比较意义，也与实际计算成本密切相关。
